\documentclass[11pt,a4paper]{article}
\usepackage[utf8]{inputenc}
\usepackage[spanish]{babel}
\usepackage{amsmath,amssymb,amsfonts}
\usepackage{geometry}
\usepackage{hyperref}
\usepackage{enumitem}
\usepackage{listings}
\usepackage{xcolor}
\usepackage{titlesec}

\geometry{margin=2.5cm}
\hypersetup{colorlinks=true, linkcolor=blue!60!black, urlcolor=blue!60!black}

\titleformat{\section}{\Large\bfseries}{Clase \thesection:}{0.5em}{}
\titleformat{\subsection}{\large\bfseries}{\thesubsection}{0.5em}{}

\lstset{
  basicstyle=\ttfamily\small,
  keywordstyle=\color{blue!70!black},
  commentstyle=\color{green!50!black},
  stringstyle=\color{red!60!black},
  breaklines=true,
  frame=single,
  backgroundcolor=\color{gray!5},
}

\title{\textbf{Inteligencia Artificial Generativa Avanzada\\y Sistemas Multi-Agente}\\[0.5em]
\large Apuntes de clase --- UTDT\\
Prof.\ Francisco Traversaro}
\author{}
\date{}

\begin{document}
\maketitle
\tableofcontents
\newpage

%=============================================================================
\section{De vectores a un LLM --- Fundamentos matemáticos}
%=============================================================================

\subsection{El camino hacia un LLM}

Todo LLM se apoya en una cadena de bloques fundamentales:

\begin{center}
\textbf{Vectores/embeddings} $\to$ \textbf{Red neuronal} $\to$ \textbf{Atenci\'on} $\to$ \textbf{Transformer} $\to$ \textbf{LLM}
\end{center}

Los s\'imbolos clave de la clase: $\mathbb{R}^n$ (espacios vectoriales), $\mathbf{w}^\top \mathbf{x}$ (producto interno), $W\mathbf{x}+\mathbf{b}$ (capa af\'in), $\phi$ (activaci\'on), $L(\theta)$ (p\'erdida), $\nabla L$ (gradiente), AD (diferenciaci\'on autom\'atica).

\subsection{Vectores, producto interno y embeddings}

\begin{itemize}
  \item Un vector $\mathbf{x} = (x_1, \ldots, x_n) \in \mathbb{R}^n$ tiene dos lecturas: $n$-tupla de coordenadas (\'algebra) y punto/direcci\'on con norma (geometr\'ia).
  \item En deep learning, todo dato (token, imagen) se codifica como vector en dimensi\'on alta: \textbf{representaci\'on distribuida}.
  \item \textbf{Embeddings}: vectores aprendidos por palabra. La geometr\'ia del espacio captura significado:
  \[
    \text{vec}(\text{rey}) - \text{vec}(\text{hombre}) + \text{vec}(\text{mujer}) \approx \text{vec}(\text{reina})
  \]
  \item \textbf{Producto interno}: $\mathbf{w}^\top \mathbf{x} = \sum_i w_i x_i = \|\mathbf{w}\| \|\mathbf{x}\| \cos\theta$. Mide alineaci\'on/proyecci\'on entre vectores. Base de la \textbf{similitud coseno}.
  \item En atenci\'on, el peso de un token sobre otro es el producto interno de sus vectores query y key (escalado por $\sqrt{d}$, normalizado con softmax).
\end{itemize}

\subsection{Aplicaciones lineales y matrices}

\begin{itemize}
  \item $W \in \mathbb{R}^{m \times n}$ representa $\mathbb{R}^n \to \mathbb{R}^m$. Cada componente de la salida es un producto interno de una fila con la entrada.
  \item Una \textbf{capa lineal} son $m$ productos internos en paralelo.
  \item Componer capas lineales = multiplicar matrices. El producto matriz--matriz domina el costo del entrenamiento $\Rightarrow$ hardware paralelo (GPU).
  \item \textbf{SVD} ($W = U\Sigma V^\top$): toda aplicaci\'on lineal es rotaci\'on $\cdot$ escala $\cdot$ rotaci\'on. Relevante para LoRA y atenci\'on de rango bajo.
\end{itemize}

\subsection{Unidad af\'in + no linealidad}

\begin{itemize}
  \item Unidad: $a = \phi(\mathbf{w}^\top \mathbf{x} + b)$. Capa: $\mathbf{h} = \phi(W\mathbf{x} + \mathbf{b})$.
  \item \textbf{Sin $\phi$, apilar capas no agrega expresividad}: la composici\'on de afines es af\'in. La no linealidad rompe eso.
  \item Activaciones comunes: ReLU, sigmoide ($\sigma$), tanh.
  \item \textbf{Teorema de aproximaci\'on universal}: una red con una capa oculta y no linealidad no polinomial aproxima cualquier funci\'on continua sobre un compacto (Cybenko 1989, Hornik 1991). No dice cu\'antas neuronas ni si el entrenamiento las encuentra.
\end{itemize}

\subsection{Aprendizaje supervisado como optimizaci\'on}

\begin{itemize}
  \item Datos $(x_i, y_i)$, modelo $f_\theta$, p\'erdida $L(\theta) = \frac{1}{N}\sum_i \ell(f_\theta(x_i), y_i)$.
  \item Entrenar = resolver $\min_\theta L(\theta)$ (en general no convexo).
  \item \textbf{Funciones de p\'erdida}:
    \begin{itemize}
      \item Regresi\'on: MSE $\ell(\hat{y}, y) = \frac{1}{2}\|\hat{y} - y\|^2$ (modelo gaussiano).
      \item Clasificaci\'on: entrop\'ia cruzada $\ell = -\sum_i y_i \log \hat{y}_i$ con $\hat{y} = \text{softmax}(z)$ (modelo categ\'orico).
    \end{itemize}
  \item Elegir una p\'erdida es, en el fondo, elegir un modelo probabilístico de los datos.
\end{itemize}

\subsection{Gradiente y descenso por gradiente}

\begin{itemize}
  \item $\nabla_\theta L \in \mathbb{R}^p$: vector de derivadas parciales, apunta a la direcci\'on de m\'aximo ascenso.
  \item Regla de actualizaci\'on: $\theta \leftarrow \theta - \eta \nabla_\theta L(\theta)$.
  \item $\eta$ grande $\to$ diverge; $\eta$ chico $\to$ converge lento.
  \item \textbf{SGD}: gradiente estimado sobre un minibatch (insesgado).
  \item \textbf{Momentum}: acumula inercia. \textbf{Adam}: momentum + escala adaptativa por coordenada. Optimizador por defecto en la pr\'actica.
\end{itemize}

\subsection{Backpropagation como diferenciaci\'on autom\'atica}

\begin{itemize}
  \item \textbf{Regla de la cadena}: $\frac{\partial z}{\partial x} = \frac{\partial z}{\partial y} \cdot \frac{\partial y}{\partial x}$.
  \item Backprop = evaluar la cadena de la salida hacia la entrada, reutilizando resultados (producto de Jacobianos).
  \item \textbf{Modo reverso}: 1 escalar de salida (la p\'erdida), $p$ par\'ametros $\Rightarrow$ gradiente completo en UNA pasada hacia atr\'as, costo $O(\text{forward})$. El modo directo costar\'ia $O(p)$.
  \item \textbf{Grafo de c\'omputo}: nodos = valores intermedios; aristas = operaciones. Forward: evaluar en orden topol\'ogico. Backward: recorrer en orden inverso.
\end{itemize}

\subsection{Notebook: Motor de autograd (micrograd)}

El laboratorio construye un motor de diferenciaci\'on autom\'atica en modo reverso en $\sim$100 l\'ineas de Python:

\begin{enumerate}
  \item \textbf{Value}: nodo del grafo. Guarda \texttt{data}, \texttt{grad} ($\partial L/\partial \text{self}$), \texttt{\_prev} (padres), \texttt{\_backward} (closure que propaga gradiente).
  \item Operaciones (\texttt{\_\_add\_\_}, \texttt{\_\_mul\_\_}, \texttt{tanh}, \texttt{exp}): cada una crea un nodo, registra padres, y define su \texttt{\_backward} con la derivada local.
  \item \textbf{backward()}: orden topol\'ogico + barrido inverso disparando cada \texttt{\_backward}.
  \item \textbf{Gradient check}: verificar anal\'itico $\approx$ num\'erico con diferencias finitas centradas. Regla de oro: \emph{nunca conf\'ies en un gradiente sin verificarlo}.
  \item \textbf{Neuron $\to$ Layer $\to$ MLP}: todo se apoya en Value. El loop de entrenamiento (forward $\to$ loss $\to$ zero\_grad $\to$ backward $\to$ update) es conceptualmente el mismo con que se entrena un LLM.
\end{enumerate}

\textbf{Punto clave}: los gradientes se \textbf{acumulan} (\texttt{+=}), no se reemplazan. Por eso \texttt{zero\_grad()} es imprescindible en cada paso.

%=============================================================================
\section{Attention is All You Need --- Transformers}
%=============================================================================

\subsection{El punto de partida}

Paper: Vaswani et al., NeurIPS 2017. Introdujo el Transformer: arquitectura basada enteramente en atenci\'on, sin recurrencia ni convoluciones. Base de los modelos de lenguaje actuales.

\subsection{Del texto a tokens}

\begin{itemize}
  \item Tres granularidades: car\'acter, subpalabra, palabra.
  \item $V$ = tama\~no del vocabulario (fijo). $T$ = longitud de la secuencia (var\'ia por texto).
  \item Trade-off: car\'acter = vocab chico pero secuencias largas; palabra = vocab enorme pero secuencias cortas. Subpalabra (BPE) busca el equilibrio.
\end{itemize}

\subsection{El bigrama: el modelo m\'as simple}

\begin{itemize}
  \item Predice el pr\'oximo car\'acter mirando \'unicamente el actual. Sin contexto, sin historia.
  \item \texttt{token\_embedding\_table} = \texttt{nn.Embedding(V, V)}: la fila $i$ son los logits del pr\'oximo token dado el token $i$.
  \item De logits a probabilidades: $\text{softmax}(\text{fila}_c)$.
  \item Loss antes de entrenar: tabla random $\to$ predicciones picudas y equivocadas $\to$ loss $> \ln(V)$ (peor que el ``ignorante honesto'' uniforme).
  \item Entrenar un bigrama = contar pares. El gradiente en la fila: $\text{grad} = p - \text{onehot}(\text{correcto})$.
  \item \textbf{Techo}: no conoce el contexto. Solo codifica ``qu\'e caracteres suelen seguir a $c$'' promediado sobre todo el corpus.
\end{itemize}

\subsection{Promediar el pasado: darle memoria}

\begin{itemize}
  \item Representaci\'on de cada token = promedio de todos los tokens hasta \'el (bag of words del pasado). Ojo causal: solo hacia atr\'as.
  \item Tres versiones equivalentes:
    \begin{enumerate}
      \item For-loop (correcto pero lento).
      \item Matriz triangular normalizada: $\text{wei} = \text{tril}/\text{sum}$, \texttt{xbow = wei @ x}.
      \item Softmax de afinidades: zeros $\to$ m\'ascara $-\infty$ $\to$ softmax $\to$ wei @ v.
    \end{enumerate}
  \item \textbf{El molde que no cambia}: (1) armar wei, (2) m\'ascara $-\infty$, (3) softmax por fila, (4) salida = wei @ v.
  \item \textbf{Problema del promedio}: descarta el \textbf{orden} y la \textbf{relevancia} (todos pesan igual).
\end{itemize}

\subsection{Self-Attention}

\begin{itemize}
  \item Cambiamos UNA pieza del molde: wei pasa a calcularse a partir de los tokens (ya no son ceros).
  \item Tres proyecciones lineales por token:
    \begin{itemize}
      \item \textbf{Query} $q = x \cdot W_Q$: ``qu\'e estoy buscando''.
      \item \textbf{Key} $k = x \cdot W_K$: ``qu\'e ofrezco para ser encontrado''.
      \item \textbf{Value} $v = x \cdot W_V$: ``qu\'e aporto si me eligen''.
    \end{itemize}
  \item Afinidad: $\text{wei}[i,j] = q_i \cdot k_j$ (producto interno).
  \item F\'ormula de una cabeza:
  \[
    \text{Attention}(Q,K,V) = \text{softmax}\!\left(\frac{QK^\top}{\sqrt{d_k}} + \text{m\'ascara}\right) \cdot V
  \]
  \item Escalado $1/\sqrt{d_k}$: control de varianza. Sin \'el, el softmax se satura y el gradiente se apaga.
  \item \textbf{Positional encoding}: la atenci\'on es ciega al orden $\to$ entrada = embedding(token) + PE(pos).
  \item ``Self'': Q, K, V salen de la misma secuencia. Si K, V vinieran de otra, ser\'ia cross-attention.
\end{itemize}

\subsection{Multi-Head Attention}

\begin{itemize}
  \item Una cabeza = un solo par $W_Q/W_K$ = una sola noci\'on de ``relevante''. Insuficiente: en lenguaje conviven varios tipos de relaci\'on (sintaxis, proximidad, correferencia).
  \item Multi-head: $H$ cabezas independientes, cada una con sus propias $W_Q, W_K, W_V$, en subespacios de tama\~no $C/H$.
  \item Concatenar + proyecci\'on $W_O$: mezcla los bloques. Entra $(B,T,C)$, sale $(B,T,C)$.
  \item \textbf{Costo}: $H$ cabezas de $C/H$ tienen los mismos par\'ametros que 1 cabeza de $C$. Expresividad pr\'acticamente gratis.
\end{itemize}

\subsection{El bloque Transformer completo}

\begin{itemize}
  \item \textbf{FFN} (Feed-Forward Network): $\text{FFN}(x) = W_2 \cdot \phi(W_1 x + b_1) + b_2$. Cómputo por token (vertical), independiente por posici\'on. Dimensi\'on interna t\'ipica $4C$.
  \item Atenci\'on \textbf{comunica} tokens; FFN \textbf{procesa} cada token.
  \item \textbf{Conexiones residuales}: $x \leftarrow x + \text{Sublayer}(x)$. Camino directo para el gradiente $\to$ combate el desvanecimiento.
  \item \textbf{LayerNorm}: normaliza cada vector (media 0, varianza 1) sobre sus $C$ componentes. Estabiliza activaciones.
  \item Bloque completo:
  \begin{align*}
    x &\leftarrow x + \text{MultiHead}(\text{LN}(x)) \\
    x &\leftarrow x + \text{FFN}(\text{LN}(x))
  \end{align*}
  \item Se apilan $N$ bloques. Capas bajas $\to$ estructura local; altas $\to$ relaciones abstractas. BERT: $N=12$, GPT-3: $N=96$.
\end{itemize}

\subsection{Encoder vs Decoder}

\begin{center}
\begin{tabular}{lll}
  \textbf{Encoder-only} & \textbf{Decoder-only} & \textbf{Encoder-decoder} \\
  BERT & GPT (lo que construimos) & T5, paper original \\
  Comprensi\'on & Generaci\'on & Secuencia a secuencia \\
\end{tabular}
\end{center}

\begin{itemize}
  \item Encoder: atenci\'on bidireccional (sin m\'ascara). Decoder: atenci\'on causal (con m\'ascara).
  \item Cross-attention: Q del decoder, K y V del encoder.
  \item Se enmascara solo cuando ver un token ser\'ia hacer trampa respecto de lo que falta generar.
\end{itemize}

%=============================================================================
\section{Detr\'as de la IA --- Tokenizers}
%=============================================================================

\subsection{Por qu\'e importa la tokenizaci\'on}

La tokenizaci\'on est\'a en el origen de muchos comportamientos an\'omalos de los LLMs. La unidad m\'inima es el \textbf{token}, no el car\'acter.

\textbf{Comportamientos que explica}:
\begin{itemize}
  \item \textbf{Deletreo}: no deletrea ni cuenta letras dentro de un token.
  \item \textbf{Aritm\'etica}: n\'umeros segmentados de formas inconsistentes.
  \item \textbf{Otros idiomas}: m\'as tokens por palabra fuera del ingl\'es $\to$ peor rendimiento y mayor costo.
  \item \textbf{Espacios}: un espacio final puede alterar la tokenizaci\'on.
  \item \textbf{Tokens raros} (SolidGoldMagikarp): token visto por el tokenizer pero casi no entrenado por el modelo $\to$ comportamiento err\'atico.
  \item \textbf{Formato}: JSON vs YAML consume distinta cantidad de tokens.
\end{itemize}

\subsection{De texto a bytes}

\begin{itemize}
  \item \textbf{Unicode}: asigna a cada car\'acter un \textbf{code point} ($\sim$150.000 definidos).
  \item Usar code points como vocabulario: $V \approx 150.000$, inestable, matrices enormes.
  \item \textbf{UTF-8}: reescribe cada code point como 1 a 4 \textbf{bytes} (0--255). Alfabeto fijo de 256, universal, sin fuera-de-vocabulario.
  \item ASCII: code point $\leq$ 127 $\to$ 1 byte (coincide con el n\'umero). Texto en ingl\'es usa menos bytes.
  \item L\'imite del vocabulario de 256: $T$ crece demasiado. Costo de atenci\'on $O(T^2)$.
\end{itemize}

\subsection{Byte-Pair Encoding (BPE)}

\begin{itemize}
  \item Mecanismo: sobre los bytes, repetir: fusionar el par adyacente m\'as frecuente en un token nuevo.
  \item Vocabulario inicial: 256 bytes (ids 0--255). Cada merge a\~nade un token: id 256, 257, \ldots
  \item Tama\~no final: $V = 256 + \text{n\'umero de merges}$ (hiperpar\'ametro de dise\~no).
  \item \textbf{El entrenamiento no tiene gradiente}: es puro conteo de frecuencias sobre un corpus.
  \item Produce dos estructuras: \textbf{merges} (par $\to$ id nuevo, en orden) y \textbf{vocab} (id $\to$ bytes).
\end{itemize}

\subsection{Encode y Decode}

\begin{itemize}
  \item \textbf{train}: descubre las reglas (por frecuencia). Se ejecuta una vez, offline.
  \item \textbf{encode(texto $\to$ ids)}: aplica los merges por \textbf{prioridad de merge} (orden en que se aprendieron), \textbf{no por frecuencia}. Texto $\to$ bytes $\to$ repetir: buscar el par presente cuyo merge se aprendi\'o primero $\to$ fusionar.
  \item \textbf{decode(ids $\to$ texto)}: lookup en vocab. id $\to$ \texttt{vocab[id]} (bytes) $\to$ concatenar $\to$ decodificar UTF-8.
  \item \textbf{Round-trip exacto}: \texttt{decode(encode(texto)) == texto} para todo texto.
  \item Error com\'un: encode por frecuencia en vez de por prioridad de merge.
\end{itemize}

\subsection{Regex splitting (GPT-2/GPT-4)}

\begin{itemize}
  \item Sin restricciones, BPE fusiona a trav\'es de l\'imites de palabra (\texttt{perro.}, \texttt{perro!} como tokens distintos).
  \item Soluci\'on: un patr\'on regex corta el texto en fragmentos (palabras, n\'umeros, puntuaci\'on, espacios) \textbf{antes} de BPE. Los merges nunca cruzan fragmentos.
  \item GPT-4 vs GPT-2: contracciones case-insensitive, n\'umeros en grupos de 1--3 d\'igitos, mejor manejo de espacios.
  \item Estos patrones se dise\~nan a mano, con sesgo hacia el ingl\'es.
\end{itemize}

\subsection{Tokens especiales}

\begin{itemize}
  \item Marcas de control fuera de BPE: \texttt{<|endoftext|>}, \texttt{<|im\_start|>}, \texttt{<|im\_end|>}, etc.
  \item Se agregan a mano con un id reservado, por encima de los bytes y merges.
  \item \texttt{encode} los separa antes de BPE. Riesgo de seguridad: hay que habilitarlos explícitamente (\texttt{allowed\_special}).
  \item ChatML: \texttt{<|im\_start|>user} ... \texttt{<|im\_end|>} \texttt{<|im\_start|>assistant}.
\end{itemize}

\subsection{Tokenizers reales y algoritmos}

\begin{center}
\begin{tabular}{llll}
  \textbf{Familia} & \textbf{Algoritmo} & \textbf{Librer\'ia} \\
  \hline
  GPT (2/3/4) & BPE byte-level & tiktoken \\
  BERT & WordPiece & BERT / HF tokenizers \\
  T5 & Unigram & SentencePiece \\
  Llama 2 / Mistral & BPE (code points) & SentencePiece \\
  Llama 3 & BPE byte-level & estilo tiktoken \\
\end{tabular}
\end{center}

\begin{itemize}
  \item \textbf{BPE}: bottom-up, fusiona pares m\'as frecuentes.
  \item \textbf{WordPiece} (BERT): bottom-up, fusiona por m\'axima verosimilitud (no por frecuencia).
  \item \textbf{Unigram}: top-down, arranca con vocabulario grande y \textbf{poda} lo que menos aporta. Usa EM para estimar probabilidades, Viterbi para segmentar.
  \item \textbf{SentencePiece}: toolkit de Google. Corre BPE o Unigram sobre code points (no bytes). \texttt{byte\_fallback} para OOV.
\end{itemize}

\subsection{Notebook: Tokenizaci\'on}

\textbf{Parte A --- Tokenizer m\'inimo desde cero}:
\begin{itemize}
  \item Implementaci\'on de \texttt{get\_stats} (contar pares) y \texttt{merge} (reemplazar par por id nuevo).
  \item 3 merges sobre \texttt{"aaabdaaabac"} $\to$ \texttt{[258, 100, 258, 97, 99]}.
  \item Sobre texto real: 20 merges $\to$ compresi\'on $\sim$1.57x.
  \item Round-trip exacto verificado.
  \item Regex splitting con patr\'on de GPT-2: cada categor\'ia (letras, n\'umeros, puntuaci\'on) queda separada.
\end{itemize}

\textbf{Parte B --- Tokenizers reales}:
\begin{itemize}
  \item tiktoken: \texttt{cl100k} (GPT-4) vs \texttt{o200k} (GPT-4o). Vocabularios crecientes (50k $\to$ 100k $\to$ 200k).
  \item AutoTokenizer: GPT-2 (BPE, prefijo \texttt{\.G}), BERT (WordPiece, prefijo \texttt{\#\#}), T5 (Unigram, símbolo \texttt{\_}).
  \item Rarezas en vivo: \texttt{strawberry} $\to$ \texttt{[st, raw, berry]}; n\'umeros partidos inconsistentemente; m\'as tokens para japon\'es que ingl\'es; JSON cuesta m\'as tokens que YAML.
\end{itemize}

\subsection{Notebook: ``\textquestiondown Qu\'e es realmente la IA?''}

Notebook introductorio usando Qwen3 8B con Ollama. Puntos clave:

\begin{itemize}
  \item Un LLM \textbf{solo predice el pr\'oximo token}. Todo lo dem\'as (memoria, personalidad, herramientas, b\'usqueda) es \textbf{software alrededor} del predictor.
  \item \textbf{raw=True}: apaga el chat template (el modelo recibe texto plano). \textbf{think=False}: apaga la cadena de razonamiento.
  \item La ``conversaci\'on'' es una ilusi\'on: un \'unico string con tokens especiales (\texttt{<|im\_start|>}, \texttt{<|im\_end|>}).
  \item \textbf{Temperatura}: controla el azar. $T=0$ $\to$ determin\'istico; $T=1$ $\to$ diverso.
  \item \textbf{Memoria}: cada llamada arranca de cero. Soluci\'on: reenviar TODO el historial.
  \item \textbf{Personalidad}: system prompt (texto al inicio del string).
  \item \textbf{Herramientas}: el modelo no calcula. Soluci\'on naive: router por keyword. Soluci\'on real: function calling.
  \item \textbf{Alucinaci\'on}: inventa con total seguridad cuando no sabe. Soluci\'on: RAG (inyectar contexto).
  \item \textbf{Alineamiento}: SFT (fine-tuning supervisado) + RLHF/preferencias. Convierte un simulador de texto en un asistente. Efecto colateral: adulaci\'on (sycophancy).
  \item Modelo base vs instruct: el base contin\'ua texto; el instruct responde a instrucciones.
\end{itemize}

%=============================================================================
\section{Embeddings}
%=============================================================================

\subsection{Parte 1: Word Embeddings}

\subsubsection{Los embeddings de un LLM se crean dentro del modelo}

La tabla de embeddings (\texttt{nn.Embedding}) es un par\'ametro m\'as, entrenado end-to-end con la misma loss de predecir el pr\'oximo token. Los modelos que vemos en esta clase son para entender el mecanismo.

\subsubsection{Similitud coseno}

\[
  \cos(\mathbf{a}, \mathbf{b}) = \frac{\mathbf{a} \cdot \mathbf{b}}{\|\mathbf{a}\| \|\mathbf{b}\|}
\]

\begin{itemize}
  \item $\cos = 1$: misma direcci\'on (0\textdegree). $\cos = 0$: ortogonales (90\textdegree). $\cos = -1$: opuestos (180\textdegree).
  \item Distancia coseno: $\text{dist} = 1 - \cos$ (m\'as chica = m\'as cerca).
  \item Compara \textbf{direcci\'on, no tama\~no}. Un vector el doble de largo pero misma direcci\'on da $\cos = 1$.
  \item Se usa coseno (no euclidiana) porque la magnitud correlaciona con frecuencia, y queremos factorizarla afuera.
\end{itemize}

\subsubsection{One-hot: identidad, cero significado}

\begin{itemize}
  \item Vector de $V$ ceros con un \'unico 1. Gigante, sparse.
  \item $\cos(\text{gato}, \text{perro}) = 0 = \cos(\text{gato}, \text{avi\'on})$: todo ortogonal a todo.
  \item \textbf{Identidad clave}: $\text{one\_hot}(i) \cdot C = C[i]$. Un embedding es leer una fila de una tabla.
\end{itemize}

\subsubsection{Hip\'otesis distribucional y conteo}

\begin{itemize}
  \item ``Una palabra se conoce por la compa\~n\'ia que mantiene'' (Firth, 1957).
  \item \textbf{Matriz de co-ocurrencia}: $M[i,j]$ = veces que $j$ cae en la ventana de $i$.
  \item Problema del conteo crudo: lo domina la frecuencia. Palabras comunes co-ocurren con todo.
  \item \textbf{PPMI} (Positive Pointwise Mutual Information): $\text{PMI}(w,c) = \log\frac{P(w,c)}{P(w) \cdot P(c)}$, recortado a 0. Saca el efecto de frecuencia.
  \item \textbf{SVD} para comprimir la matriz. Cada palabra pasa de $|$contextos$|$ dimensiones a $d$.
  \item Emerge geometr\'ia (temas se agrupan), pero es ruidosa y no escala.
\end{itemize}

\subsubsection{Word2Vec: aprender prediciendo}

\begin{itemize}
  \item \textbf{Truco central}: inventar una tarea de predicci\'on descartable. La tarea es la excusa; los vectores son el premio.
  \item \textbf{Skip-gram}: dada una palabra central, predecir sus vecinas dentro de una ventana. Self-supervised.
  \item Score: $\text{score}(\text{centro}, \text{contexto}) = \mathbf{v}_{\text{centro}} \cdot \mathbf{u}_{\text{contexto}}$.
  \item \textbf{Cuello de botella}: softmax sobre todo $V$ ($O(|V|)$ por par).
  \item \textbf{Negative sampling}: de ``elegir entre $V$'' a ``\textquestiondown real o falso?''. $k+1$ problemas binarios con sigmoide. $O(k)$ en vez de $O(V)$.
  \item Resultado: analog\'ias como aritm\'etica: $\text{rey} - \text{hombre} + \text{mujer} \approx \text{reina}$.
  \item Descubrimiento (Levy \& Goldberg 2014): skip-gram con negative sampling factoriza impl\'icitamente una matriz de PMI.
\end{itemize}

\subsubsection{GloVe: contar y despu\'es aprender}

\begin{itemize}
  \item Combina lo mejor de contar (global, eficiente) y predecir (patrones ricos).
  \item El objeto: el \textbf{cociente} $P(k|\text{ice})/P(k|\text{steam})$ discrimina; la probabilidad sola, no.
  \item Modelo: $\mathbf{v}_i \cdot \mathbf{u}_k + b_i + b_k = \log X_{ik}$. Regresi\'on ponderada sobre la matriz de conteos.
  \item \textbf{No es una red}: profundidad cero, sin no linealidad, sin softmax. Modelo bilineal.
  \item Los tres (SVD+PPMI, skip-gram, GloVe) factorizan la misma matriz de asociaci\'on.
  \item Las analog\'ias no son emergentes: son la restricci\'on de dise\~no.
\end{itemize}

\subsubsection{Techo: static embeddings y polisemia}

Un vector fijo por palabra, ciego al contexto. ``Bank'' tiene dos sentidos pero un solo vector. La resoluci\'on: embeddings \textbf{contextuales} (self-attention / Transformer).

\textbf{Arco completo}: one-hot (identidad) $\to$ static (significado, ciego al contexto) $\to$ contextual (Transformer).

\subsubsection{Sesgos}

Los embeddings heredan los sesgos del corpus. Ejemplo: $\text{man}:\text{doctor}::\text{woman}:\text{nurse}$. Medir y mitigar estos sesgos es parte de responsible AI.

\subsection{Parte 2: Sentence Embeddings}

\subsubsection{El problema}

Las aplicaciones operan sobre \textbf{textos}, no palabras sueltas. Necesitamos un vector por oraci\'on, de largo fijo, que capture el significado completo.

\subsubsection{Baseline: promediar word vectors}

\begin{itemize}
  \item $\text{emb}(\text{oraci\'on}) = \frac{1}{n}\sum_{i=1}^{n} \mathbf{v}_i$. Centro de masa sem\'antico.
  \item Funciona sorprendentemente bien como piso. Pero tiene tres agujeros:
    \begin{enumerate}
      \item \textbf{Orden}: ``the dog bit the man'' y ``the man bit the dog'' $\to$ mismo vector ($\cos = 1$).
      \item \textbf{Diluci\'on}: oraciones largas promedian mucho relleno.
      \item \textbf{OOV}: palabras fuera de vocabulario se ignoran.
    \end{enumerate}
\end{itemize}

\subsubsection{BERT: el encoder}

\begin{itemize}
  \item Encoder-only, bidireccional. Entrenado con masked language modeling.
  \item Produce un vector contextual por token. No genera texto.
  \item \textbf{BERT crudo no sirve para coseno}: nunca lo entrenaron para que cerca en significado $\Rightarrow$ cerca en coseno. Pooleado, rinde peor que promediar GloVe en STS.
\end{itemize}

\subsubsection{SBERT (Sentence-BERT)}

\begin{itemize}
  \item SBERT = BERT + pooling + entrenamiento \textbf{contrastivo}.
  \item \textbf{Pooling}: de un vector por token a uno por oraci\'on. Mean pooling (m\'as com\'un) o token [CLS].
  \item \textbf{Entrenamiento contrastivo} (siamés): el mismo encoder embebe dos oraciones; la loss acerca las parecidas y aleja las distintas. Optimiza la geometr\'ia del coseno.
  \item Datos: pares etiquetados (par\'afrasis, duplicados, NLI). Negativos gratis: in-batch negatives.
  \item Cierra los tres agujeros: self-attention lee el orden, pondera importancia, tokeniza en subwords (sin OOV).
  \item Interfaz: \texttt{model.encode("una oración")} $\to$ un vector.
\end{itemize}

\subsubsection{Embedders basados en decoder (GPT)}

\begin{itemize}
  \item Misma receta: vectores contextuales $\to$ poolear $\to$ fine-tuning contrastivo.
  \item Diferencia: se lee el \textbf{\'ultimo} token (el \'unico que vio toda la oraci\'on), no el [CLS].
  \item No se empieza de cero: arranca del GPT pre-entrenado. El pre-entrenamiento pone el \textbf{conocimiento}; el contrastivo pone la \textbf{m\'etrica}.
\end{itemize}

\subsubsection{Tres ``embeddings'' distintos}

\begin{enumerate}
  \item \textbf{De entrada}: vector fijo por token del vocabulario (lookup).
  \item \textbf{Contextual por token}: salida de las capas con atenci\'on.
  \item \textbf{Del texto} (el de RAG): UN solo vector para toda la query/documento.
\end{enumerate}

\subsubsection{B\'usqueda sem\'antica = la ``R'' de RAG}

\begin{itemize}
  \item Query + documentos $\to$ coseno $\to$ top-$k$. Eso es la recuperaci\'on de RAG.
  \item Match por \textbf{significado}, no por palabras. Una query puede no compartir palabras con el documento correcto y recuperarlo.
  \item \textbf{Invariante cr\'itica}: query y documentos deben embeberse con el \textbf{mismo modelo}. El coseno solo tiene sentido dentro de un mismo espacio vectorial. Cambiar de embedder $\to$ re-indexar todo.
  \item Sim\'etrico (comparar frases similares) vs asim\'etrico (query corta $\to$ pasaje largo).
\end{itemize}

\subsubsection{Los tres modelos de un RAG}

\begin{enumerate}
  \item \textbf{Embedder} que indexa la base y embebe la query (mismo modelo).
  \item \textbf{LLM generador} que escribe la respuesta (recibe texto, no vectores).
  \item Se comunican por texto, no por vectores.
\end{enumerate}

\subsubsection{Qu\'e embedder usar (2026)}

\begin{itemize}
  \item Benchmark: \textbf{MTEB}. Criterios: idioma, tama\~no/latencia, licencia.
  \item APIs pagas: OpenAI, Cohere, Google (fuertes pero lock-in).
  \item Open source competitivo: Qwen3-Embedding, BGE-M3 (multiling\"ue).
  \item Livianos y locales: nomic-embed-text (137M, CPU), bge-m3.
  \item Regla: benchmarkear en tus propios datos.
\end{itemize}

\subsubsection{Lo que falta para RAG en serio (Semana 8)}

\begin{itemize}
  \item \textbf{Chunking}: partir documentos largos en pedazos indexables.
  \item \textbf{Vector store}: Chroma, pgvector, Qdrant (en vez de lista de numpy).
  \item \textbf{Reranking}: segundo paso que afina el top-$k$.
\end{itemize}

\subsection{Notebook: Word Embeddings (N1)}

Recorrido ejecutable completo:
\begin{itemize}
  \item \S1: Coseno con n\'umeros y plot (vectores 2D).
  \item \S2: One-hot $\to$ lookup: \texttt{one\_hot @ C == C[i]} verificado con \texttt{torch.allclose}.
  \item \S3: Co-ocurrencia + PPMI + SVD desde cero sobre Reuters. Heatmaps de la matriz. Emerge geometr\'ia por temas (petr\'oleo vs granos).
  \item \S4: Word2Vec real con gensim sobre text8 ($\sim$17M palabras). Vecinos de ``king'' y analog\'ia $\text{king} - \text{man} + \text{woman} \approx \text{queen}$.
  \item \S5: GloVe pre-entrenado (400k palabras). Analog\'ias limpias. Polisemia de ``bank'' (dominada por sentido financiero). Sesgos: $\text{doctor} - \text{man} + \text{woman} \approx \text{nurse}$.
  \item \S7: Skip-gram + negative sampling desde cero en PyTorch. Producto punto $\to$ logsigmoid. Con 50k tokens los vecinos son ruidosos: la geometr\'ia necesita escala.
\end{itemize}

\subsection{Notebook: Sentence Embeddings (N2)}

\begin{itemize}
  \item \S2: Baseline promedio. $\cos(\text{king}, \text{monarch}) = 0.929$ (alto), $\cos(\text{king}, \text{pizza}) = 0.759$ (bajo). Agujero del orden: dog/man $\to \cos = 1.0$.
  \item \S3: SBERT (\texttt{all-MiniLM-L6-v2} y multilingüe). Dog/man $\to \cos = 0.98$ (ya no idénticos).
  \item \S4: Corpus de conceptos del curso. Heatmap de similitud: word2vec--negative sampling se iluminan juntos.
  \item \S5: B\'usqueda sem\'antica a mano y con \texttt{util.semantic\_search}. Query sin palabras en com\'un recupera el documento correcto.
  \item \S5: Invariante: \texttt{cos(qA, qB) = 0.236} con dos modelos distintos $\to$ espacios incompatibles.
  \item \S6: RAG m\'inimo Ollama-nativo. Embedder: nomic-embed-text (768-dim). Generador: qwen3:8b. Sin RAG inventa fechas; con RAG responde desde la base.
\end{itemize}

%=============================================================================
\section{Pretraining y Fine-tuning}
%=============================================================================

\subsection{Parte 1: Preentrenamiento --- De la web cruda al modelo base}

\subsubsection{Tres etapas: masa, silueta, cara}

\begin{itemize}
  \item El preentrenamiento (unsupervised) construye una ``masa'' gigantesca de patrones. El SFT le da silueta. RLHF es la capa final que lo vuelve asistente.
  \item Casi todo el c\'omputo y los datos se gastan en el preentrenamiento. El ajuste posterior es, en comparaci\'on, una capa fina.
\end{itemize}

\subsubsection{Los datos: curaci\'on emp\'irica}

\textbf{Perplexity}:
\begin{itemize}
  \item $\text{PPL} = \exp\!\bigl(-\frac{1}{N}\sum_i \log P(w_i|\text{contexto})\bigr) = \exp(\text{cross-entropy})$.
  \item Interpretaci\'on: n\'umero efectivo de opciones equiprobables entre las que el modelo elige en cada paso. PPL = 1 (perfecto), PPL = $|V|$ (in\'util).
  \item \textbf{Limitaci\'on}: mide fluidez con un estilo de texto, no competencia real. No compara entre tokenizers.
\end{itemize}

\textbf{Ablations}:
\begin{itemize}
  \item Experimento controlado: cambiar un solo factor (filtro, fuente, dedup) y dejar todo lo dem\'as fijo.
  \item Se eval\'ua con tareas con respuesta conocida (accuracy), no con perplexity.
  \item Modelos chicos y baratos, confiando en que el ranking relativo se mantiene a escala.
  \item \textbf{Regla}: cada decisi\'on de curaci\'on se gana con un experimento, no con intuici\'on.
\end{itemize}

\textbf{Pipeline de curaci\'on} (embudo de 8 pasos):
\begin{enumerate}
  \item \textbf{URL}: lista negra de dominios (porno, spam, malware).
  \item \textbf{Extracci\'on de texto}: quitar HTML, men\'us, ads. Solo el art\'iculo.
  \item \textbf{Idioma}: detector + ruteo por idioma. Upsampling de lenguas con poca data.
  \item \textbf{Gopher}: heur\'isticas de calidad (textos cortos, exceso de s\'imbolos, repetici\'on).
  \item \textbf{MinHash}: deduplicaci\'on de casi-copias. Deduplicar de m\'as puede ser contraproducente.
  \item \textbf{Filtros C4}: sin puntuaci\'on, lorem ipsum, boilerplate.
  \item \textbf{Filtros propios}: reglas ad-hoc para basura espec\'ifica del corpus.
  \item \textbf{PII}: reemplazar emails, IPs, tel\'efonos por placeholders.
\end{enumerate}

Referencia abierta: FineWeb / FineWeb-2 (HuggingFace). Cada compa\~n\'ia tiene su ``libro de recetas'' cerrado, pero el esquema es com\'un.

\textbf{Lecci\'on central}: calidad $>$ cantidad. Filtrar mejor rinde m\'as que agregar m\'as data cruda.

\subsubsection{La escala: Scaling Laws}

\begin{itemize}
  \item Tres perillas atadas: $C \approx 6 \cdot N \cdot D$ (c\'omputo $\approx$ 6 $\times$ par\'ametros $\times$ tokens).
  \item A $C$ fijo, $N$ y $D$ se compensan: m\'as par\'ametros $\Rightarrow$ menos tokens, y viceversa.
  \item \textbf{Power law}: la loss baja de forma predecible (recta en log-log) al escalar $N$, $D$ o $C$, hasta una \emph{irreducible loss} (entrop\'ia intr\'inseca del lenguaje).
  \item \textbf{Predicci\'on del run}: entrenar modelos chicos, fitear la curva, extrapolar. GPT-4 predijo su loss final desde runs 1000$\times$ m\'as chicos.
  \item \textbf{Chinchilla} (DeepMind, 2022): para entrenamiento compute-\'optimo, $\sim$20 tokens por par\'ametro. Chinchilla (70B, 1.4T tokens) le gan\'o a Gopher (280B, 300B tokens) con el mismo c\'omputo.
  \item \textbf{Giro moderno}: sobre-entrenar por inferencia. Llama 3 8B: $\sim$1900 tokens/par\'ametro. Modelo chico, caro una vez, barato de servir.
\end{itemize}

\subsubsection{El modelo base}

\begin{itemize}
  \item Producto del pretraining: un ``simulador de documentos de internet''. Competencia amplia, comportamiento crudo.
  \item \textbf{No se comporta como asistente}: ante ``\textquestiondown Cu\'al es la capital de Francia?'' puede continuar con m\'as preguntas (completa el patr\'on en vez de responder).
  \item El conocimiento ya est\'a adentro: con few-shot, el base responde bien. La capacidad estaba; falta activarla.
  \item \textbf{Competencia} (lo que sabe) vs \textbf{comportamiento} (c\'omo se presenta): dos cosas distintas. Pretraining compra competencia (caro, una vez). Comportamiento es una capa final barata.
\end{itemize}

\subsubsection{Continued pretraining vs fine-tuning}

\begin{center}
\begin{tabular}{lll}
  & \textbf{Continued PT} & \textbf{Fine-tuning} \\
  \hline
  Para qu\'e & Conocimiento (dominio/idioma) & Comportamiento (responder, formato) \\
  Objetivo & Next-token (self-supervised) & Pares supervisados \\
  Datos & Texto plano, mucho & Pares, poco \\
  Costo & US\$ cientos de miles & Chico \\
\end{tabular}
\end{center}

\textbf{Regla}: conocimiento nuevo $\to$ continu\'a pre-entrenando; comportamiento nuevo $\to$ fine-tuning. No met\'es un dominio entero con un dataset chico de instrucciones.

Ejemplo: modelo base + continued pretraining sobre corpus coreano ($\sim$200B tokens, $\sim$US\$0.2M) = modelo coreano. Mezclar coreano + ingl\'es para evitar catastrophic forgetting.

\subsection{Parte 2: Fine-tuning eficiente}

\subsubsection{Instruction tuning (SFT)}

\begin{itemize}
  \item Entrenar sobre pares (instrucci\'on $\to$ respuesta ideal): \textbf{demostraciones} de buen comportamiento.
  \item Generaliza porque la variedad de tareas (resumir, traducir, explicar, clasificar) ense\~na la \textbf{habilidad general} de seguir instrucciones, no un cat\'alogo de respuestas.
  \item \textbf{LIMA} (Less Is More for Alignment): mil ejemplos buenos y variados le ganan a un set enorme y ruidoso. Calidad $>$ cantidad.
  \item \textbf{Masking del prompt}: la loss se calcula solo sobre los tokens de la respuesta. La instrucci\'on condiciona; la respuesta ense\~na.
  \item \textbf{Chat template}: serializa roles (system/user/assistant) con tokens especiales (ChatML: \texttt{<|im\_start|>}, \texttt{<|im\_end|>}). Mismo template en entrenamiento e inferencia.
  \item SFT full fine-tuning + reward model + PPO = pipeline est\'andar de la industria (InstructGPT, 2022).
\end{itemize}

\subsubsection{PEFT: Parameter-Efficient Fine-Tuning}

\begin{itemize}
  \item Full fine-tuning mueve todos los pesos. A escala real, no entra en una GPU: $\sim$16 bytes/par\'ametro (pesos + gradiente + estados de Adam). Un 7B $\to$ $\sim$112 GB.
  \item PEFT congela el base y entrena $<1\%$ de par\'ametros. El conocimiento ya est\'a; para cambiar comportamiento alcanza un ajuste chico.
  \item Tres estrategias:
    \begin{itemize}
      \item \textbf{Adapters}: capas chicas intercaladas en el Transformer (cuello de botella $d \to r \to d$). Suman latencia.
      \item \textbf{Soft prompts / Prefix tuning}: vectores aprendidos en la entrada (o en todas las capas de atenci\'on). Gastan contexto.
      \item \textbf{LoRA}: ajuste de bajo rango sobre las matrices existentes. Se fusiona $\to$ cero overhead en inferencia.
    \end{itemize}
\end{itemize}

\subsubsection{LoRA (Low-Rank Adaptation)}

\begin{itemize}
  \item Hallazgo emp\'irico: el $\Delta W$ del fine-tuning tiene \textbf{rango bajo}.
  \item En vez de aprender $\Delta W \in \mathbb{R}^{d \times d}$, aprender $B \cdot A$ con $A \in \mathbb{R}^{r \times d}$, $B \in \mathbb{R}^{d \times r}$, $r \ll d$.
  \item Forward: $h = W x + \frac{\alpha}{r} B A x$. $W$ congelado, solo $A$ y $B$ entrenables.
  \item Inicializaci\'on: $A$ al azar, $B$ en ceros $\Rightarrow$ $\Delta W = 0$ al inicio (modelo arranca id\'entico al base).
  \item \textbf{La cuenta}: una capa $4096 \times 4096$ = 16.7M par\'ametros. Con LoRA $r=8$: $2 \times 8 \times 4096 = 65.5$K $\approx 0.4\%$.
  \item Hiperpar\'ametros: $r$ (8--16), $\alpha$ (16--32), target modules (q, k, v, o de atenci\'on).
  \item Beneficios: memoria m\'inima para entrenar + adapters livianos ($\sim$MB) e intercambiables.
\end{itemize}

\subsubsection{QLoRA}

\begin{itemize}
  \item LoRA resolvi\'o lo entrenable, pero el base congelado sigue ocupando memoria ($\sim$14 GB en fp16 para un 7B).
  \item QLoRA = LoRA + base cuantizado a \textbf{4 bits}. Almacenamiento a 4 bits, c\'omputo de-cuantizado al vuelo a 16 bits.
  \item \textbf{NF4} (4-bit NormalFloat): 16 niveles en los cuantiles de una normal. M\'as valores donde hay m\'as pesos $\to$ menos error de redondeo.
  \item \textbf{Double quantization}: cuantiza tambi\'en las constantes de escala.
  \item \textbf{Paged optimizers}: ante picos de memoria, descarga estados del optimizador a CPU RAM.
  \item Resultado: fine-tuning de un 65B en una sola GPU de 48 GB. Un 7B entra en una T4 de 16 GB (Colab gratis).
  \item Casi sin costo de calidad: 4-bit $\approx$ 16-bit en pr\'actica.
\end{itemize}

\subsubsection{S\'intesis}

\begin{itemize}
  \item \textbf{Una sola loss} (next-token) de punta a punta: preentrenamiento $\to$ SFT $\to$ LoRA $\to$ QLoRA.
  \item Los datos cambiaron una vez: de texto crudo a demostraciones instrucci\'on $\to$ respuesta.
  \item En cada escal\'on se mueven menos pesos: full FT (100\%) $\to$ LoRA ($<1\%$) $\to$ QLoRA ($<1\%$ + base 4-bit).
  \item \textbf{Techo del SFT}: clona demostraciones, llega tan lejos como sus ejemplos. No captura preferencias. Siguiente paso: RLHF / DPO.
\end{itemize}

\subsection{Notebook: Fine-tuning eficiente}

Arco completo sobre SmolLM2-1.7B (base) con T4 gratis:

\begin{enumerate}
  \item \textbf{Cargar base en 4-bit}: \texttt{BitsAndBytesConfig} con NF4 + double quant. 1.71B pesos, $\sim$1.0 GB (vs 3.42 GB en fp16).
  \item \textbf{Base no se comporta}: ante ``Explain photosynthesis'' genera m\'as preguntas en loop; ante ``Capital of France'' responde pero no frena.
  \item \textbf{Chat template}: ChatML (\texttt{<|im\_start|>}/\texttt{<|im\_end|>}) definido manualmente (el base no lo trae, pero los tokens ya existen en el vocabulario).
  \item \textbf{Datos}: Dolly 15k, subset de 1000. 8 categor\'ias (open\_qa, classification, summarization\ldots). Formateados con chat template.
  \item \textbf{LoRA a mano}: \texttt{LoRALayer} (18 l\'ineas): $A$ al azar, $B$ en ceros. \texttt{LinearWithLoRA}: $Wx + (\alpha/r) BAx$. Verificaci\'on: salida id\'entica al inicio ($\Delta W = 0$).
  \item \textbf{Con \texttt{peft}}: \texttt{LoraConfig(r=8, alpha=16, targets=[q,k,v,o])}. Trainable: 3.1M / 1.71B = \textbf{0.18\%}. Adapters en fp32, base en 4-bit.
  \item \textbf{Entrenar}: 200 pasos con \texttt{SFTTrainer}, LR $2 \times 10^{-4}$. Loss baja de $\sim$3.5 a $\sim$2.3.
  \item \textbf{Antes/despu\'es}: ahora responde (``Paris'') y frena. Bugs cr\'iticos: \texttt{model.eval()} (apagar dropout de LoRA) y \texttt{eos\_token\_id} (frenar en \texttt{<|im\_end|>}).
  \item \textbf{Adapter guardado}: 6.1 MB (vs 3.42 GB del modelo completo).
\end{enumerate}

%=============================================================================
\section{Alineamiento y Evaluaci\'on}
%=============================================================================

\subsection{Parte 1: Alineamiento}

\subsubsection{La analog\'ia de Karpathy}

Tres formas de estudiar un libro = tres etapas de entrenamiento:
\begin{enumerate}
  \item \textbf{Pretraining} = exposici\'on, conocimiento de fondo.
  \item \textbf{SFT} = problemas resueltos, imitaci\'on.
  \item \textbf{RL} = problemas de pr\'actica, probar y corregir.
\end{enumerate}

\subsubsection{Dominios verificables: el loop de RL}

\begin{itemize}
  \item Donde hay respuesta conocida (matem\'atica, c\'odigo), el modelo genera muchas soluciones, un \textbf{verificador} filtra las correctas, y se refuerzan.
  \item Una estrategia confiable (llega m\'as seguido a la respuesta) se refuerza m\'as, sin que nadie la etiquete.
  \item \textbf{El razonamiento emerge}: al optimizar por acertar, las respuestas se alargan solas (el modelo descubre que ``pensar en voz alta'' ayuda). Es lo que produce R1, o1, Qwen3 think.
  \item \textbf{Formato \texttt{<think>}}: el contenido del razonamiento emerge del RL; el tag \texttt{<think>} se impone (template o cold-start SFT).
  \item RL supera al que imita: como solo importa llegar a la respuesta, el modelo encuentra caminos que ning\'un humano usar\'ia (AlphaGo, ``movimiento 37'').
\end{itemize}

\subsubsection{Dominios no verificables: preferencias}

\begin{itemize}
  \item ``Un chiste sobre pel\'icanos'', ``resum\'i bien esto'': no hay respuesta correcta contra qu\'e chequear.
  \item \textbf{Comparar, no puntuar}: mostrar al humano varias respuestas y pedir que las ordene. Ordenar es m\'as f\'acil que puntuar. La unidad m\'inima es un par (chosen, rejected).
\end{itemize}

\subsubsection{Reward Model}

\begin{itemize}
  \item Un juez aprendido que reemplaza al verificador. Recibe (prompt, respuesta) y devuelve $r \in \mathbb{R}$.
  \item Nace de un GPT con la cabeza cambiada: mismas capas de Transformer, distinta cabeza (escalar en vez de distribuci\'on sobre vocabulario).
  \item Se entrena con Bradley-Terry: $\min -\log \sigma(r(\text{chosen}) - r(\text{rejected}))$.
  \item Un ranking de $K$ respuestas se desarma en $\binom{K}{2}$ pares.
\end{itemize}

\subsubsection{RLHF}

\begin{itemize}
  \item ``Pol\'itica'' = el LLM que entren\'as. La pol\'itica genera, el reward model punt\'ua, se empuja hacia puntajes altos.
  \item \textbf{Reward hacking}: el RL encuentra agujeros del reward model (entradas sin sentido con puntaje alt\'isimo, e.g.\ ``the the the'' $\to$ 1.0).
  \item \textbf{Freno KL}: $\max \mathbb{E}[r(y)] - \beta \cdot \text{KL}(\pi \| \pi_{\text{ref}})$. Se congela una copia del SFT como ancla. $\beta$ = largo de la correa.
  \item RLHF es un ajuste fino acotado: se corta a los pocos cientos de pasos (el reward es hackeable). Pipeline caro: hasta 4 modelos simult\'aneos.
\end{itemize}

\subsubsection{DPO (Direct Preference Optimization)}

\begin{itemize}
  \item Mismo objetivo que RLHF (subir lo preferido sin despegarse del SFT), pero \textbf{sin reward model} y \textbf{sin PPO}: una sola loss sobre los pares.
  \item Explota que el juez y el generador son el mismo Transformer con distinta cabeza $\to$ colapsa las dos redes en una.
  \item PPO y DPO est\'an al mismo nivel: dos caminos desde los mismos pares.
\end{itemize}

\subsubsection{Constitutional AI (RLAIF)}

\begin{itemize}
  \item Reemplaza al humano por principios en lenguaje natural (``constituci\'on'') + un modelo que juzga.
  \item \textbf{Fase 1}: criticar y revisar. El modelo genera, se critica contra un principio, y reescribe mejor $\to$ dataset para SFT.
  \item \textbf{Fase 2}: preferencias por IA. Un modelo compara dos respuestas seg\'un la constituci\'on $\to$ pares (chosen, rejected) sin humano.
  \item El humano subi\'o un nivel: de etiquetar casos a escribir principios. Es la receta de Anthropic (Claude).
  \item Misma trampa: el juez-IA es hackeable.
\end{itemize}

\subsubsection{S\'intesis del alineamiento}

\begin{center}
\begin{tabular}{ll}
  \textbf{\textquestiondown Qui\'en punt\'ua?} & \textbf{Consecuencia} \\
  \hline
  Verificador (objetivo) & Corre indefinido, razonamiento emerge, supera humanos \\
  Reward model (proxy) & Hackeable, correa KL, techo acotado \\
\end{tabular}
\end{center}

\subsection{Parte 2: Evaluaci\'on}

\subsubsection{Benchmarks}

\begin{itemize}
  \item Bater\'ia fija de tareas con respuesta conocida $\to$ un n\'umero comparable. MMLU, HellaSwag, GSM8K, HumanEval.
  \item \textbf{Mecanismo}: no es generaci\'on; se rankean log-probs de las opciones. Mide reconocimiento, no producci\'on.
  \item \textbf{Saturaci\'on}: los modelos de frontera superan el benchmark $\to$ la comunidad responde con sucesores m\'as dif\'iciles (MMLU $\to$ MMLU-Pro $\to$ GPQA Diamond $\to$ HLE).
  \item \textbf{Contaminaci\'on}: benchmarks online pueden haber entrado en los datos de entrenamiento. \textquestiondown Razon\'o o record\'o?
  \item Miden capacidad, no preferencia. 90 en MMLU e insufrible de usar.
\end{itemize}

\subsubsection{Evaluaci\'on por preferencia}

\begin{itemize}
  \item El mismo par (chosen/rejected) que entren\'o al modelo, ahora como instrumento de medici\'on.
  \item \textbf{LLM-as-a-judge}: un modelo fuerte elige entre dos respuestas. Coincide con humanos $\sim$80\%.
  \item Sesgos del juez: posici\'on (prefiere A o B seg\'un orden), \textbf{verbosidad} (premia la m\'as larga), auto-preferencia (favorece su propio estilo).
  \item \textbf{Chatbot Arena}: batallas ciegas entre modelos, votadas por humanos, ranking Elo.
  \item \textbf{Unificaci\'on}: Elo = Bradley-Terry = la misma $\sigma(\text{rating}_A - \text{rating}_B)$ que entrena el reward model y que usa DPO.
\end{itemize}

\subsubsection{Modos de falla}

\textbf{Hallucinations}:
\begin{itemize}
  \item El modelo produce afirmaciones falsas con confianza. No es bug: el pretraining premia fluidez, no verdad. RLHF puede empeorar la calibraci\'on (sycophancy).
  \item \textbf{Detecci\'on}: SelfCheckGPT --- muestrear $N$ veces, medir desacuerdo. Un hecho conocido coincide; una invenci\'on diverge.
  \item \textbf{Mitigaci\'on}: entrenar el ``no s\'e'', tool use / web search, RAG.
\end{itemize}

\textbf{Sesgos}:
\begin{itemize}
  \item Cambiar un \'unico t\'ermino (nombre, g\'enero, nacionalidad) y medir si la respuesta cambia sistem\'aticamente.
  \item A escala: StereoSet (estereotipos), ToxiGen (toxicidad impl\'icita).
\end{itemize}

\textbf{Goodhart / Reward hacking}:
\begin{itemize}
  \item ``Cuando una medida se convierte en objetivo, deja de ser una buena medida.''
  \item Es exactamente el reward hacking del entrenamiento, ahora en evaluaci\'on. Ejemplo: optimizar para el juez LLM = hacer respuestas m\'as largas. Win-rate sube, calidad real no.
  \item La \'unica evaluaci\'on que Goodhart no rompe: la verificable (donde hay verdad chequeable).
\end{itemize}

\subsubsection{Responsible AI}

\begin{itemize}
  \item Taxonom\'ia de Weidinger et al.\ (2022): 21 riesgos en 6 \'areas (discriminaci\'on/hate speech, information hazards, misinformation, usos maliciosos, da\~nos de la interacci\'on, ambientales/socioecon\'omicos).
  \item Seis elementos: fairness, inclusiveness, reliability \& safety, privacy \& security, transparency, accountability.
  \item \textbf{Guardrail}: un segundo pase que pregunta ``\textquestiondown esto es seguro?'' antes de devolver la respuesta. Es Constitutional AI en inferencia.
  \item La misma constituci\'on que en entrenamiento defin\'ia ``qu\'e es bueno'' reaparece como guardrail que define ``qu\'e es seguro mostrar''.
\end{itemize}

\subsubsection{S\'intesis de evaluaci\'on}

\begin{itemize}
  \item Un solo objeto ---la comparaci\'on de preferencias--- entren\'o al modelo y lo mide. Alineamiento y evaluaci\'on son dos caras de ``comparar respuestas''.
  \item El mismo riesgo aparece en ambos: reward hacking = Goodhart.
  \item Tensi\'on de frontera: a medida que los modelos superan la capacidad humana de juzgar (scalable oversight), evaluar se vuelve el problema dif\'icil.
  \item \textbf{Puente a lo que sigue}: el modelo alineado y evaluado tiene una debilidad pr\'actica: inventa. La soluci\'on: RAG (darle una fuente).
\end{itemize}

%=============================================================================
\newpage
\section*{Lecturas recomendadas}
\addcontentsline{toc}{section}{Lecturas recomendadas}

\begin{description}
  \item[Clase 1] Goodfellow, Bengio \& Courville, \emph{Deep Learning} (cap.\ 2, 4, 6.5). CS231n (Stanford). Repo micrograd (A.\ Karpathy).
  \item[Clase 2] Vaswani et al.\ (2017), \emph{Attention Is All You Need}. Karpathy, nanoGPT. Brown et al.\ (2020), \emph{Language Models are Few-Shot Learners} (GPT-3).
  \item[Clase 3] Sennrich et al.\ (2015), BPE en NLP. Radford et al.\ (2019), GPT-2 \S2.2. Video: Karpathy, \emph{Let's build the GPT Tokenizer}. Repo minbpe.
  \item[Clase 4] Lena Voita, \emph{Word Embeddings}. CS224N (Stanford). Jay Alammar, \emph{The Illustrated Word2vec}. Reimers \& Gurevych (2019), \emph{Sentence-BERT}. Leaderboard MTEB.
  \item[Clase 5] Hu et al.\ (2021), \emph{LoRA}. Dettmers et al.\ (2023), \emph{QLoRA}. Zhou et al.\ (2023), \emph{LIMA}. Ouyang et al.\ (2022), \emph{InstructGPT}. FineWeb (HuggingFace). Chinchilla (Hoffmann et al., 2022).
  \item[Clase 6] Karpathy, \emph{Deep Dive into LLMs}. Ouyang et al.\ (2022), InstructGPT. Rafailov et al.\ (2023), DPO. DeepSeek-R1 (2025). Bai et al.\ (2022), Constitutional AI. Zheng et al.\ (2023), LLM-as-a-Judge. Weidinger et al.\ (2022), taxonom\'ia de riesgos.
\end{description}

\end{document}
